\documentclass[10pt,a4paper]{amsart}
\usepackage[margin=19mm]{geometry}
\usepackage{amsmath,amssymb,mathtools}
\usepackage[scr=rsfs,cal=euler]{mathalfa}
\usepackage{xfrac}
\usepackage[unicode,hidelinks,bookmarks=false]{hyperref}
\newcommand{\ie}{i.e.\ }
\newcommand{\cf}{cf.\ }
\newcommand{\resp}{respectively\ }
\newcommand{\Subsection}{\subsection}
\providecommand{\nobreakdash}{\nobreak\hbox{-}}
\setlength{\emergencystretch}{2em}
\multlinegap0pt
\usepackage[shortlabels]{enumitem}
\usepackage{bm}
\usepackage{amssymb}
\usepackage{xcolor}
\usepackage{cancel}
\newtheorem{lemma}{Lemma}
\newtheorem{proposition}{Proposition}
\usepackage{cleveref}
\crefname{lemma}{Lemma}{Lemmas}
\crefname{proposition}{Proposition}{Propositions}
\newcommand{\bfGamma}{{\boldsymbol \Gamma}}
\newcommand{\bfX}{{\bf X}}
\newcommand{\bfZ}{{\bf Z}}
\newcommand{\bfbeta}{{\boldsymbol \beta}}
\newcommand{\bfgamma}{{\boldsymbol \gamma}}
\newcommand{\bfx}{{\bf x}}
\newcommand{\bfy}{ {\bf y}}
\newcommand{\bfz}{{\bf z}}
\def\eqref#1{equation~\ref{#1}}
\newcommand{\tr}[1]{\textcolor{red}{#1}}
\title{Stability of Transformers under Layer Normalization}
\author{Kelvin Kan, Xingjian Li, Benjamin J. Zhang, Tuhin Sahai, Stanley Osher, Krishna Kumar, Markos A. Katsoulakis}
\date{Source: arXiv:2510.09904v2 (CC BY 4.0)}
\begin{document}
\maketitle

\noindent\emph{Source and licence.} Stability of Transformers under Layer Normalization. Version arXiv:2510.09904v2 (CC BY 4.0), distributed by arXiv under the Creative Commons Attribution 4.0 International licence, http://creativecommons.org/licenses/by/4.0/, as recorded in the arXiv licence metadata for this exact version and shown on the abstract page https://arxiv.org/abs/2510.09904v2. Full licence notice: \texttt{licenses/arxiv-2510.09904-CC-BY-4.0.txt}.\par\smallskip

\noindent\textit{Editorial scope.} Counted unit: the article's proposition prop:Bellman (source0.tex lines 1423--1429). The article's complete original proof of it is delivered with it (source0.tex lines 1431--1500, one \texttt{proof} environment of 70 lines). No mathematics is written, repaired or re-derived: the statement and the proof are the article's own lines, in the article's own notation, and no retained line is substituted or re-flowed.  Retained uncounted as the source-local context the proof needs: lines 791--797; lines 1397--1408; lines 1418--1420.  Nothing was omitted from any retained span, so the package carries no omission record.  No retained line contains a citation, so the wrapper carries no bibliography and no bibliography entry.  No retained line contains a figure environment, an image-inclusion command or drawing code, so no image is delivered and the package declares no asset dependency.  Editorial formatting only, in addition: an AMS \texttt{amsart} preamble that restates the article's own declarations and macros used by the retained lines, namely the environments \texttt{lemma}, \texttt{proposition} on separate counters, together with the standard packages those lines need.  The retained environments print in this document as Lemma 1, Proposition 1 in order of appearance, whereas the article prints the same items with the numbering of its own full document.  The complete native source of the article as distributed in the arXiv e-print for this exact version is bundled unchanged as \texttt{sources/186841/source0.tex}.
\begin{lemma}\label{lemma:ellipsoid}
    The layer normalization output $\bfz=\text{LN}(\bfx; \bfgamma, \bfbeta)$ lies on the ellipsoid 
    $$
    \mathcal{E} = \left\{ \mathbf{z} \in \mathbb{R}^d : (\mathbf{z} - \bfbeta)^\top \bfGamma^{-2} (\mathbf{z} - \bfbeta) = d \right\}, 
    $$
    where $\bfGamma = {\rm diag}(\bfgamma)\in \mathbb{R}^{d\times d}$.
\end{lemma}

\subsection{Assumptions}
In our derivations, we use the following assumptions.
\begin{enumerate}[label=(A\arabic*)]
    \item \label{assump:one}The functions $L_i: \mathbb{R}^{d\times n} \to \mathbb{R} \cup \{\infty\}$ are convex, proper, and lower-semicontinuous (l.s.c.),
    \item \label{assump:two}$L_i$'s are coercive, that is,
    \[
    \frac{L_i(f)}{\| f\|_F} \to \infty \quad \text{as $\| f\|_F^2 \to \infty$ }
    \]
    \item \label{assump:three}$L_i$'s are convex and satisfies $L_i(\bfX) \geq 0$ for all $i$ and $\bfX \in \mathbb{R}^{d \times n}$, and 
    \item \label{assump:four}the terminal cost $G(\cdot, \bfy)$ is continuous and bounded below: there exists $0< M < \infty$ such that $G(\cdot, \bfy) \geq -M$ for all $\bfX \in \mathbb{R}^{d \times n}$.
\end{enumerate}
We remark that for Pre-LN, $L_i$'s are the constant zero function, and for Peri-LN, $L_i$'s are an indicator function over an ellipsoid (\Cref{lemma:ellipsoid}), which is a strictly convex set. In these two cases, the first assumption is satisfied. 

\begin{equation}\label{eq:value_function}
    \Phi_{i, \bfy}(\bfX) = \inf_{\{ f_j\}_{j \geq i} } \left\{ G(\bfX_D, \bfy) + \Delta t \sum_{j \geq i} L_j(f_j(\bfX_j)) \; \middle| \; \bfX_i=\bfX \right\},
\end{equation}

\begin{proposition}[Bellman Equation]\label{prop:Bellman}
    Under assumptions~\ref{assump:one}-\ref{assump:four}, the value function~\eqref{eq:value_function} satisfies the Bellman equation
    \begin{equation}\label{eq:value_function_recursive}
        \Phi_{i, \bfy}(\bfX) = \inf_{f} \left\{ \Phi_{i+\frac{1}{2}, \bfy}(\bfX + \Delta t \cdot f) + \Delta t \cdot L_i(f(\bfX)) \right\},
    \end{equation}
    for $i=0,\frac{1}{2},1,...,D-\frac{1}{2}$. Moreover, the infimum is attained: for every $\bfX \in \mathbb{R}^{d \times n}$ and every $i$, there exists an optimal velocity achieving the infimum.
\end{proposition}

\begin{proof}
    Fix $i \in \{ 0,\frac{1}{2},1,...,D-\frac{1}{2} \}$, and recall that we have the given terminal condition $\Phi_{D, \bfy}(\bfX) = G(\bfX, \bfy)$. Denote the right-hand-side of~\eqref{eq:value_function_recursive} by
    \begin{equation}\label{eq:value_function_recursive_psi}
        \tilde{\Phi}_{i, \bfy}(\bfX) := \inf_f \Psi_{\bfX, \bfy}(f), \quad \Psi_{\bfX, \bfy}(f) :=  \Phi_{i+\frac{1}{2}, \bfy}(\bfX + \Delta t \cdot f) + \Delta t \cdot L_i(f(\bfX)) .
    \end{equation}
    Our goal is to show that $\tilde{\Phi}_{i, \bfy}(\bfX) = {\Phi}_{i, \bfy}(\bfX)$ and that the infimum is attained.

    \textbf{Step 1} (${\Phi}_{i, \bfy}(\bfX) \geq \tilde{\Phi}_{i, \bfy}(\bfX)$)\textbf{.} Let $\bfX=\bfX_i$ and
    $\{ f_j\}_{j \geq i}$ be any control sequence.  Denote $f:=f_i$ as the first control. The total cost from~\eqref{eq:value_function} can be split as follows
    \begin{equation}
        \Delta t \sum_{j \geq i} L_j(f_j(\bfX_j)) + G(\bfX_D, \bfy)  = \Delta t \cdot L_i(f(\bfX)) + \underbrace{\Delta t \sum_{j \geq i+\frac{1}{2}} L_j(f_j(\bfX_j)) + G(\bfX_D, \bfy)}_{\text{tail cost from $\bfX_{i+\frac{1}{2}}$}}. 
    \end{equation}
    By~\eqref{eq:value_function} and the definition of infimum, the tail cost from $\bfX_{i+\frac{1}{2}}$ is at least $\Phi_{i+\frac{1}{2}, \bfy}(\bfX + \Delta t \cdot f)$. Thus we have
    \[
    \Delta t \sum_{j \geq i} L_j(f_j(\bfX_j)) + G(\bfX_D, \bfy) \geq \Delta t \cdot L_i(f(\bfX)) + \Phi_{i+\frac{1}{2}, \bfy}(\bfX + \Delta t \cdot f) = \Psi_{\bfX, \bfy}(f) \geq \tilde{\Phi}_{i, \bfy}(\bfX) .
    \]
    Taking infimum over all control sequence yields ${\Phi}_{i, \bfy}(\bfX) \geq \tilde{\Phi}_{i, \bfy}(\bfX)$.

    \textbf{Step 2} (${\Phi}_{i, \bfy}(\bfX) \leq \tilde{\Phi}_{i, \bfy}(\bfX)$)\textbf{.} Fix any $f \in \mathbb{R}^{d \times n}$ and $\epsilon > 0$. Set $\bfZ=\bfX + \Delta t \cdot f$. By definition of infimum and~\eqref{eq:value_function}, there exists a tail control sequence $\{ f_{j}^\epsilon \}_{j \geq i+\frac{1}{2}}$ starting from state $\bfX_{i+\frac{1}{2}}=\bfZ$ with tail cost at most $\Phi_{i+\frac{1}{2}, \bfy}(\bfZ) + \epsilon$. Consider a full control sequence from $(i, \bfX)$ which uses $f$ at layer $i$ and the tail control sequence $\{ f_{j}^\epsilon \}_{j \geq i+\frac{1}{2}}$ for all subsequent layers. This yields a hidden state trajectory with $\bfX_i=\bfX$, $\bfX_{i+\frac{1}{2}}=\bfZ$ and total cost
    \begin{equation}
        \Delta t \cdot L_i(f(\bfX)) + [\text{tail cost from $\bfZ$}] \leq \Delta t \cdot L_i(f(\bfX)) + \Phi_{i+\frac{1}{2}, \bfy}(\bfZ) +\epsilon = \Psi_{\bfX, \bfy}(f) + \epsilon.
    \end{equation}
    We remark that the constructed control sequence starts at $\bfX_i=\bfX$ and layer $i$, and the value function~\eqref{eq:value_function} is defined as the infimum cost, also starting at $\bfX_i=\bfX$ and layer $i$, thus we have
    \begin{equation}
        \Phi_{i, \bfy}(\bfX) \leq \Psi_{\bfX, \bfy}(f) + \epsilon.
    \end{equation}
    Since $\epsilon$ is arbitrary, taking $\epsilon \to 0$ yields
    \begin{equation}
        \Phi_{i, \bfy}(\bfX) \leq \Psi_{\bfX, \bfy}(f)
    \end{equation}
    Morevoer, taking infimum over the constructed sequence yields
    \begin{equation}
        \Phi_{i, \bfy}(\bfX) \leq \inf_f\Psi_{\bfX, \bfy}(f) =  \tilde{\Phi}_{i, \bfy}(\bfX).
    \end{equation}

    \textbf{Step 3 (Attainment of the infimum).} We show that the infimum in the Bellman equation
    \[
    {\Phi}_{i, \bfy}(\bfX) = \inf_{f} \left\{ \Phi_{i+\frac{1}{2}, \bfy}(\bfX + \Delta t \cdot f) + \Delta t \cdot L_i(f(\bfX)) \right\}
    \]
    is attained for some $f$. The arguments rely on the coercivity of $L_i$ (Assumption~\ref{assump:two}).

    Our argument uses a backward inductive argument starting with $i=D-\frac{1}{2}$, by the terminal condition $\Phi_{D, \bfy}(\bfX) = G(\bfX, \bfy)$, we have
    \begin{equation}\label{eq:Bellman_terminal}
    {\Phi}_{D-\frac{1}{2}, \bfy}(\bfX) = \inf_{f} \left\{ G(\bfX + \Delta t \cdot f, \bfy) + \Delta t \cdot L_{D-\frac{1}{2}}(f(\bfX)) \right\}.
    \end{equation}
    In the objective function, the first term $G$ is continuous and lower bounded by $-M>-\infty$ (assumption~\ref{assump:four}), and the second term $\Delta t \cdot L_{D-\frac{1}{2}}$ is ls.c., coercive, and lowered bounded by 0 (assumptions \ref{assump:one}, \ref{assump:two}, and \ref{assump:three}). Thus, the objective $ f \mapsto G(\bfX + \Delta t \cdot f, \bfy) + \Delta t \cdot L_{D-\frac{1}{2}}(f(\bfX))$ is coercive, bounded below and l.s.c.. This implies that, for any $c>0$, the sublevel set 
    \[
    \left\{ f\middle| G(\bfX + \Delta t \cdot f, \bfy) + \Delta t \cdot L_{D-\frac{1}{2}}(f(\bfX)) \leq c \right\}
    \]
    is bounded (coercivity) and closed (by l.s.c.) and hence compact. Thus, the Weierstrass extreme value theorem~\tr{cite} implies that the infimum is attained at some $f$. 

    In addition, the infimum objective in~\eqref{eq:Bellman_terminal} is jointly l.s.c., and the infimum of $f$ is taken over a compact set (due to the coercivity of $L$), ${\Phi}_{D-\frac{1}{2}, \bfy}(\bfX)$ is l.s.c.. Moreover, ${\Phi}_{D-\frac{1}{2}, \bfy}(\bfX) = \inf_{f} \{ G(\bfX + \Delta t \cdot f, \bfy) + \Delta t \cdot L_{D-\frac{1}{2}}(f(\bfX)) \} \geq 0-M=-M$, so it is bounded below. Thus, we can repeat the same arguments to iteratively show that the infimum is attained at some $f$ for ${\Phi}_{i, \bfy}(\bfX) $ for $i=D-1, D-\frac{3}{2},...,0$. 

    We remark that, we see that in the proof of attainment of infimum, a key assumption is the coercivity of $L_i$'s. However, recall that under Pre-LN, $L_i's$ are the constant zero function since no restriction is imposed onto the velocity field (output of the layers). Thus, under Pre-LN, the coercivity assumption does not hold.

    In more details, in the Pre-LN training problem, we only have the terminal cost $G$ and no running cost $L_i$'s to regulate the  velocity $f_i$'s. Thus, there are infinitely many velocity $f_i$'s that can attain the infimum of the Bellman equation~\eqref{eq:value_function_recursive}, including highly irregular ones. For instance, since $L_i$'s are not coercive, the velocity can have unbounded magnitude. This is the fundamental reason of the ill-posedness of the Pre-LN training problem.
    % We first establish that 
    % \[
    % \Psi_{\bfX, \bfy}(f) =  \Phi_{i+\frac{1}{2}, \bfy}(\bfX + \Delta t \cdot f) + \Delta t \cdot L_i(f(\bfX))
    % \]
    % is l.s.c. in $f$ and bounded below. This is because the second term $\Delta t \cdot L_i(f(\bfX))$ is l.s.c. by assumption~\ref{assump:two}, and that the first term $\Phi_{i+\frac{1}{2}, \bfy}(\bfX + \Delta t \cdot f)$ is also l.s.c. and bounded below. The lower-semicontinuity and boundedness of the first term is obtained by an backward inductive argument. First, the terminal $\Phi_{D, \bfy}=G(\cdot, \bfy)$ is l.s.c. and bounded below by some $-M<0$ by assumption~\ref{assump:four}. Next, for all $0\leq i<D$, since $L_i \geq 0$ by assumption~\ref{assump:three}, 
    % \[\Phi_{i, \bfy}(\bfX) = \inf_f \Psi_{\bfX, \bfy}(f) = \inf_f \{ \Phi_{i+\frac{1}{2}, \bfy}(\bfX + \Delta t \cdot f) + \Delta t \cdot L_i(f(\bfX)) \} \geq 0 + -M=-M.\]
    % This induces the lower boundedness of $\Psi_{\bfX, \bfy}$ as follows
    % \[
    % \Psi_{\bfX, \bfy}(f) = \Phi_{i+\frac{1}{2}, \bfy}(\bfX + \Delta t \cdot f) + \Delta t \cdot L_i(f(\bfX)) 
    % \geq -M + \Delta t \cdot L_i(f(\bfX)).
    % \]
    % Moreover, the recursive definition of the value function~\eqref{eq:value_function_recursive} also implies that $\Phi_{i, \bfy}$ is l.s.c. for all $i$.
    % The coercivity of $L_i$ (assumption~\ref{assump:two})
\end{proof}

\end{document}
